\chapter{Literature Review}
\label{chapterlabel2}

\section{The Urban Night and Night-time Activity}

The night has long been treated in urban research as an appendage of daytime activity, with its own temporal and spatial logic receiving relatively little independent attention \citep{VanLiempt2015}. Yet the urban night is not simply a continuation of the day, nor a transitional phase in which urban functions gradually wind down.

Employment, leisure, socialising, care, maintenance and return journeys all continue after conventional working hours, but at different times, in different places and under different conditions.

These differences produce patterns of urban life that are distinct from daytime commuting rhythms \citep{Schwanen2012,VanLiempt2015}.

Research that treats the night as a distinct urban period has developed largely around the concept of the ``night-time economy'' (NTE). Early work focused particularly on catering, bars, cultural consumption and entertainment, highlighting the contribution of night-time activity to urban economic development and cultural vitality.

However, viewing the night primarily through leisure and consumption can overlook its wider social dimensions. \citet{Shaw2010,Shaw2014} and \citet{Hadfield2015} show that the urban night also involves labour, governance, regulation, safety, exclusion and inequality, while different groups experience unequal conditions of access to and participation in night-time urban space \citep{Plyushteva2019}. The urban night is thus not a homogeneous period defined solely by nightlife and leisure, but a setting in which different activities, populations and social conditions coexist.

This diversity is also reflected in the spatial organisation of night-time activity. Compared with the relatively stable rhythms that daytime activity forms around employment centres, educational facilities and major commuting corridors, night-time activities may be distributed across central entertainment districts, transport hubs, residential neighbourhoods, healthcare facilities, logistics areas and concentrations of night-time employment. Opening hours, public service availability, working arrangements and spatial regulation together shape when and where these activities occur. The same location may serve different urban functions by day and by night, while different parts of the city may exhibit different patterns of activity timing, duration and directional flow \citep{Schwanen2012,Kim2020,Shaw2022}.

\section{Night-time Transportation and Mobility}

The diversity of night-time activities is accompanied by a different set of travel conditions after dark. Different groups use the same transport system at different times and for different purposes \citep{Shaw2022,Palm2023}. Understanding night-time mobility requires attention both to the conditions under which travel takes place and to the temporal relationship between activities and transport services.

First, the operating conditions of night-time travel differ from those during the day. Using a discrete choice model, \citet{Kapitza2022} finds that lighting, service availability and perceived safety influence mode choice at night. Other studies similarly show that personal safety can affect travellers' willingness to use public transport. Research in Brussels and Recife \citep{Farina2022}, together with analysis of the night-bus network in Sofia \citep{Plyushteva2020}, demonstrates that night-time mobility depends not only on the existence of transport services but also on the conditions under which those services can be used. Together, these factors influence both mode choice and the practical usability of public transport at night.

Another challenge for night-time mobility is the mismatch between the temporal structure of urban activity and transport provision. Operating hours, service frequency, interchange conditions and information availability affect whether destinations can be reached at the required time. Two locations may be well connected within the transport network, but that connection becomes less useful if the required service has been reduced or withdrawn by the time it is needed \citep{McArthur2019,Ryan2023,Pan2025}. This temporal mismatch affects a wide range of night-time travellers. Research on non-standard working hours shows that altered schedules can expose workers to transport conditions different from those experienced during the day \citep{Plyushteva2019,Lim2024,Kapitza2025}, but similar constraints apply to anyone whose activity extends beyond the hours of full-service provision, including those returning from leisure, social or care activities.

This complexity is also compounded by the fact that night-time activity can be difficult to identify within datasets dominated by daytime patterns. \citet{Kim2020} demonstrates this point using functional principal component analysis on hourly pedestrian traffic in Seoul: night-time vitality patterns emerge only when extracted as a separate component from the overwhelming daytime trend, and their temporal boundaries vary across districts according to local urban functions. The implication for transport research is that night-time usage patterns are unlikely to be adequately captured as a residual component of all-day analysis and instead require dedicated examination.

\section{Smart-card Data, Temporal Profiles and Socio-spatial Interpretation}

Studying these variations requires data that capture both when and where public transport activity occurs. Smart-card and operational transport data provide large-scale, continuous records of passenger transactions. Such records usually do not directly reveal passengers' travel purposes or their wider activity contexts. However, their scale, continuity and temporal resolution make it possible to examine how ridership varies within the day, across the week and between locations \citep{Zhong2016,Manley2018,Cats2024}.

Smart-card research has developed at both individual and location-based scales. Individual-level studies compare passengers' transactions across dates and time periods to identify regular or recurring travel behaviour \citep{Zhong2016,Manley2018}. Location-based studies instead aggregate ridership at stations, routes or areas and examine how patterns vary within the day and across the week \citep{Gan2020,Zhang2021,Cheng2024}. In this second approach, stations and areas provide spatial units through which differences in urban activity rhythms can be observed.

As research has moved beyond describing total ridership towards identifying recurring temporal structures, clustering and other unsupervised classification approaches have become widely used. \citet{Briand2016}, for example, use a mixture model to identify latent temporal usage types in public transport data. \citet{ElMahrsi2017} combine smart-card data processing and clustering to summarise large-scale urban mobility records, while \citet{Cheng2024} use fuzzy clustering to examine within-day and within-week ridership patterns at London rail stations. Other studies have used approaches including k-means and hierarchical clustering \citep{Zhao2020}. Despite differences in method, this literature demonstrates that complex temporal ridership records can be organised into a smaller number of interpretable profiles that share similar patterns of use.

Temporal profiling can also support the interpretation of transport patterns in relation to urban structure. Existing studies have linked ridership profiles with land use, urban functions and spatial characteristics to understand the different roles played by stations or areas within wider urban activity systems. \citet{Gan2020}, for example, use metro-station ridership profiles to examine urban mobility patterns, while \citet{Zhang2021} use smart-card data to study changes in London's urban structure. Other research combines transport records with household travel surveys, socio-demographic information or points-of-interest data to investigate associations between mobility patterns and factors such as jobs-housing structure, population characteristics and likely activity purposes \citep{Long2015,Zhang2019,SariAslam2021,Yang2023}.

This approach is particularly relevant to night-time public transport. Differences in night-time activity may appear in the timing of evening ridership, the persistence of activity after midnight, changes in directional flow and differences between weekdays and weekends. These patterns occur within areas that also differ in employment, residential, land-use and socio-demographic characteristics. \citet{McArthur2019} and \citet{Smeds2020} examine London's night-time transport from temporal and socio-spatial perspectives, while the London Night Workers Classification developed by \citet{Mavrogeni2025} provides a framework for comparing the geography of night-time workers across London. Together, these studies provide a basis for interpreting temporal public transport patterns in relation to the spatial structure of London's night-time activities and workforce.

\section{Synthesis: Research Gap and Thesis Positioning}

From current review, three linked strands of research frame this dissertation. Research on the urban night shows that night-time mobility is produced by multiple forms of activity and by transport conditions that vary across time and space \citep{Schwanen2012,VanLiempt2015,Shaw2022}. Smart-card research demonstrates that temporal transport records can be organised into recurring usage profiles \citep{Briand2016,ElMahrsi2017,Gan2020}. A related body of work shows that such profiles can be interpreted alongside land-use, employment and socio-spatial information. Taken together, these strands establish a pathway from night-time transport records to temporal typology and external area contextualisation, but they also reveal three connected gaps.

First, although research on the urban night has examined labour, accessibility, service provision and transport equity, it has less often identified recurring types of observed public-transport use within a dedicated night-time analytical window. Smart-card studies, by contrast, have mainly examined all-day travel, routine commuting, peak periods or general station functions, leaving night-time activity embedded within the wider daily cycle.

Second, comparative attention to the temporal, directional and spatial organisation of night-time Rail and Bus use remains limited. Because the two modes differ in their data structures and operational characteristics, this gap calls for separate mode-specific analyses followed by cautious descriptive synthesis rather than a pooled or formally equivalent comparison.

Third, existing applications provide limited evidence on how transport-derived night-time typologies correspond to an independently constructed geography of night work and to wider socio-spatial conditions. This limits the interpretation of where different observed-use patterns occur and how they are embedded within the urban night.

The dissertation addresses the three gaps in sequence: it identifies night-time usage types separately for Rail stations and bus-served LSOAs; characterises their temporal, directional, day-type and spatial structures and their descriptive cross-mode differences; and examines their associations with the London Night Workers Classification and broader area context. These contextual associations are used to interpret the environments in which the patterns occur, not to infer individual passenger identities, travel purposes, causal mechanisms, unmet demand or service deficiency.

The next chapter describes the data and methods used to implement this research framework.
